\documentclass[12pt]{article}

% 1 inch margin per side
\usepackage{fullpage}
% Includes graphics
\usepackage{graphicx}
% Packages for Mathematics, by the American Mathematical Society
\usepackage{amssymb}
\usepackage{amsmath}
% Deactivate hyphenation
\usepackage[none]{hyphenat}
% Jump between paragraphs - no indents
\usepackage{parskip}
% UTF-8 and spanish for this template
\usepackage[spanish]{babel}
\usepackage[utf8]{inputenc}
% Links of document
\usepackage{hyperref}
\usepackage{fancyhdr}
\setlength{\headheight}{15.2pt}
\setlength{\headsep}{5pt}
\pagestyle{fancy}

\newcommand{\N}{\mathbb{N}}
\newcommand{\Exp}[1]{\mathcal{E}_{#1}}
\newcommand{\List}[1]{\mathcal{L}_{#1}}
\newcommand{\EN}{\Exp{\N}}
\newcommand{\LN}{\List{\N}}

\newcommand{\comment}[1]{}
\newcommand{\lb}{\\~\\}
\newcommand{\eop}{_{\square}}
\newcommand{\hsig}{\hat{\sigma}}
\newcommand{\ra}{\rightarrow}
\newcommand{\lra}{\leftrightarrow}

% Here you can change my name for yours
\newcommand{\alumno}{Martín Cadagan Guzmán}
% Here you can change "Tarea 1" for your assignment
\rhead{Tarea 1 - \alumno}

\begin{document}
\thispagestyle{empty}
\begin{minipage}{2.3cm}
	

% Insert here the top left corner img like the logo of your university
% Remember to have a folder with all your images
\includegraphics[width=2cm]{img/logo.pdf}
\vspace{0.5cm} 
\begin{minipage}{\linewidth}
\textsc{\raggedright \footnotesize
Pontificia Universidad Católica de Chile \\
% Insert here your faculty/ Dept of studies
Departamento de Ciencias de la Computación \\
% Insert here you course
IIC2143 - Diseño detallado de software \\}
\end{minipage}


% Title
\begin{center}
\vspace{0.5cm}
{\huge\bf Tarea 1}\\
\vspace{0.2cm}
\today\\
\vspace{0.2cm}


\footnotesize{2º semestre 2020 }\\
\vspace{0.2cm}
\footnotesize{\alumno}
\rule{\textwidth}{0.05mm}
\end{center}

% Start your awnsers from here :)
\section*{Respuestas}

\subsection*{Pregunta 1}
\textbf{a)}

% Template for centered figures
\begin{figure}[h!]
	\centering
	\includegraphics[scale=0.5]{img/logo.pdf}
	\caption{You can change this text for a figure description}
\end{figure}

% Sample text
Suponga que D es cerrado y acotado. Si $f$ presenta discontinuidades en el conjunto D, entonces el problema podría igualmente tener soluciones. ¿Cierto o no? ¿Es Chefctio un impostor o no?

\end{document}
